\section{Ascending and descending conditions}
\label{am-ch6-chains}
\begin{proposition}[Maximal condition]
\label{am-prop-6-1}
\dcref{6.1}
\source{atiyah-macdonald-commutative-algebra:page-0083}
For a partially ordered set $L$, every increasing sequence is eventually
stationary iff every nonempty subset has a maximal element.
\end{proposition}
\begin{proof}
If a nonempty subset has no maximal element, choose successively a strictly
increasing sequence.  Conversely, the set of terms of an increasing sequence
has a maximal term, after which the sequence is constant.
\end{proof}
\begin{definition}[Noetherian and Artinian modules]
\label{am-def-noeth-artin-module}
\source{atiyah-macdonald-commutative-algebra:page-0083}
An $A$-module is Noetherian if its submodules satisfy the ascending chain
condition (a.c.c.), equivalently the maximal condition; it is Artinian if they
satisfy the descending chain condition (d.c.c.), equivalently the minimal
condition.  A ring is Noetherian or Artinian when its ideals have the respective
property.
\end{definition}
\begin{proposition}[Noetherian iff every submodule is finite]
\label{am-prop-6-2}
\dcref{6.2}
\source{atiyah-macdonald-commutative-algebra:page-0084}
An $A$-module is Noetherian iff every submodule is finitely generated.
\end{proposition}
\begin{proof}
If $N$ is a submodule, the finitely generated submodules of $N$ have a maximal
member; adjoining an element outside it contradicts maximality, so it equals
$N$.  Conversely, the union of an ascending chain is a submodule and is
finitely generated; all its generators occur in one term, so the chain stops.
\end{proof}
\begin{proposition}[Chain conditions and short exact sequences]
\label{am-prop-6-3}
\dcref{6.3}
\source{atiyah-macdonald-commutative-algebra:page-0084}
In a short exact sequence $0\to M'\to M\to M''\to0$, $M$ is Noetherian iff
both $M'$ and $M''$ are Noetherian, and $M$ is Artinian iff both are Artinian.
\end{proposition}
\begin{proof}
Chains in a submodule or quotient of $M$ inherit stationarity.  Conversely,
for a chain $L_i$ in $M$, its inverse images in $M'$ and images in $M''$ both
stabilize; then $L_i$ stabilizes as well.  Reverse inclusions for Artinian.
\end{proof}
\begin{corollary}[Finite direct sums]
\label{am-cor-6-4}
\dcref{6.4}
\source{atiyah-macdonald-commutative-algebra:page-0085}
Finite direct sums of Noetherian (respectively Artinian) modules are Noetherian
(respectively Artinian).
\uses{am-prop-6-3}
\end{corollary}
\begin{proof}
Induct on the number of summands using the coordinate short exact sequence.
\end{proof}
\begin{proposition}[Finite modules over chain-condition rings]
\label{am-prop-6-5}
\dcref{6.5}
\source{atiyah-macdonald-commutative-algebra:page-0085}
If $A$ is Noetherian (respectively Artinian), every finitely generated
$A$-module is Noetherian (respectively Artinian).
\end{proposition}
\begin{proof}
A finite module is a quotient of $A^n$ (Proposition~\ref{am-prop-2-3}); apply
Corollary~\ref{am-cor-6-4} and Proposition~\ref{am-prop-6-3}.
\uses{am-prop-2-3,am-cor-6-4,am-prop-6-3}
\end{proof}
\begin{proposition}[Quotients of chain-condition rings]
\label{am-prop-6-6}
\dcref{6.6}
\source{atiyah-macdonald-commutative-algebra:page-0085}
If $A$ is Noetherian (respectively Artinian), then $A/\aideal$ is so for every
ideal $\aideal$.
\end{proposition}
\begin{proof}
The quotient is a quotient module of $A$ and its $A$-submodules are exactly its
$A/\aideal$-submodules.
\end{proof}
\begin{definition}[Length and composition series]
\label{am-def-composition-series}
\source{atiyah-macdonald-commutative-algebra:page-0085}
A chain $M=M_0\supsetneq\cdots\supsetneq M_n=0$ has length $n$.  It is a
composition series when every factor $M_{i-1}/M_i$ is simple.
\end{definition}
\begin{proposition}[Jordan--Hölder length]
\label{am-prop-6-7}
\dcref{6.7}
\source{atiyah-macdonald-commutative-algebra:page-0086}
If $M$ has a composition series of length $n$, every composition series has
length $n$, and every chain extends to one.
\end{proposition}
\begin{proof}
Let $\ell(M)$ be the least composition length.  Intersect a shortest series with
a submodule $N$; the induced factors are simple or equal, giving
$\ell(N)\le\ell(M)$, with equality only when $N=M$.  Any chain strictly lowers
this number, so has length at most $\ell(M)$.  A composition series has length
at least and therefore exactly $\ell(M)$; inserting terms into a nonmaximal
chain until this bound is reached gives a composition series.
\end{proof}
\begin{proposition}[Finite length criterion]
\label{am-prop-6-8}
\dcref{6.8}
\source{atiyah-macdonald-commutative-algebra:page-0086}
An $A$-module has a composition series iff it is both Noetherian and Artinian.
\end{proposition}
\begin{proof}
A composition series bounds the length of every chain.  Conversely, repeatedly
choose maximal proper submodules using the maximal condition; the descending
chain terminates by the minimal condition, and each factor is simple.
\end{proof}
\begin{definition}[Finite length]
\label{am-def-finite-length}
\source{atiyah-macdonald-commutative-algebra:page-0086}
When these conditions hold, the common composition length is denoted $\length(M)$;
Jordan--Hölder identifies the multisets of simple factors of any two series.
\end{definition}
\begin{proposition}[Additivity of length]
\label{am-prop-6-9}
\dcref{6.9}
\source{atiyah-macdonald-commutative-algebra:page-0086}
For a short exact sequence of finite-length modules,
$\length(M)=\length(M')+\length(M'')$.
\end{proposition}
\begin{proof}
Splice a composition series of $M'$ with the inverse images in $M$ of a
composition series of $M''$; the resulting series has the claimed length.
\end{proof}
\begin{proposition}[Vector-space chain conditions]
\label{am-prop-6-10}
\dcref{6.10}
\source{atiyah-macdonald-commutative-algebra:page-0087}
For a vector space $V$ over a field, finite dimension, finite length, a.c.c.,
and d.c.c. are equivalent, and length equals dimension.
\end{proposition}
\begin{proof}
Finite-dimensional spaces have the usual finite flags.  If $V$ is infinite,
choose infinitely many independent vectors; their successive spans give a
strictly increasing chain, while the spans of the tails give a strictly
decreasing chain.
\end{proof}
\begin{corollary}[A product of maximal ideals]
\label{am-cor-6-11}
\dcref{6.11}
\source{atiyah-macdonald-commutative-algebra:page-0087}
If $(0)=\m_1\cdots\m_n$ is a product of maximal ideals, then $A$ is Noetherian
iff Artinian.
\end{corollary}
\begin{proof}
Use the finite chain $A\supset\m_1\supset\m_1\m_2\supset\cdots\supset0$.
Each factor is a vector space over a residue field, so its a.c.c. and d.c.c.
are equivalent by Proposition~\ref{am-prop-6-10}; Proposition~\ref{am-prop-6-3}
transfers this equivalence to $A$.
\uses{am-prop-6-10,am-prop-6-3}
\end{proof}
